\documentclass[11pt]{article}

\usepackage{amsmath,amssymb,amstext,amsthm}
\usepackage{geometry}
\usepackage{fancyhdr}
\usepackage[pdftex]{graphicx}
\usepackage{xcolor}

\geometry{
  verbose,
  dvips,
  width=430pt, marginparsep=0pt, marginparwidth=0pt,
  top=90pt, headheight=12pt, headsep=20pt, footskip=30pt, bottom=80pt}

\setlength{\parskip}{\medskipamount}
\setlength{\parindent}{0pt}

\linespread{1.05}

\newtheorem{theorem}{Theorem}
\newtheorem{lemma}[theorem]{Lemma}
\newtheorem{proposition}[theorem]{Proposition}
\newtheorem{corollary}[theorem]{Corollary}
\newtheorem{conjecture}[theorem]{Conjecture}
\newtheorem{obs}{Observation}
\newtheorem{clm}{Claim}
\newtheorem{ques}{Question}
\newtheorem{problem}{Problem}

\newtheorem{ex}{Example}


\newcommand{\spc}{\hspace{0.08in}}
\newcommand{\vs}{\vspace*{.1in}}
\newcommand{\E}{\mathbb{E}}
\newcommand{\Prob}{\mathbb{P}}
\newcommand{\edit}[1]{\emph{\textcolor{blue}{#1}}}


\renewenvironment{proof}{{\noindent \textbf{\textit{Proof.}}}}
{\hfill $\Box$ \vspace*{0.1in}}
\newenvironment{proofc}{{\noindent \textit{Proof of Claim.}}}
{\hfill $\Box$}


\begin{document}

\begin{LARGE}\begin{center}\bf 14.4 Tangent Planes and Normal Lines\end{center}
\end{LARGE}

\vs\vs



\section{Warm Up}
\textbf{Problem 1:} Find the plane going through $P(-1,1,2)$ normal to $2\hat{i}+3\hat{j}-\hat{k}$.

\textbf{Problem 2:} Finds the line going through $P(-1,1,2)$ in the direction of $2\hat{i}+3\hat{j}-\hat{k}$.



\section{Motivation}

Now that we know how to take the derivative at ANY direction not just $x$ and $y$, we can ask other interesting questions in 3D that we don't have in 2D. For example we can now have a tangent \emph{plane} to a surface at a point. Here we utilize what we have learned to find the equation for a tangent plane along with a normal line!

\section{Definitions \& Theorems}

\begin{enumerate}
\item In order to find the tangent plane to a point on a surface $z = f(x,y)$ it boils down to finding the normal vector to the plane. We know that any constant plane $z= c$ will intersect the surface at a level curve which we project onto the $xy$ plane (2-dimensional).

\item Since all continuous curves can be represented by a vector function, we get that out level curve can be written as,

\begin{equation*}
    \vec{r}(t) = x(t)\hat{i} + y(t) \hat{j}
\end{equation*}

\item For any level curve at $z=c$, we see that $f(x(t),y(t)) = c$. Taking the derivative with respect to $t$ of both sides, we get, $f'(x(t),y(t)) = 0$ and then using the chain rule we see,

\begin{align*}
    \frac{\partial f}{\partial x}\frac{dx}{dt} + \frac{\partial f}{\partial y}\frac{dy}{dt} &= 0 \\
   \left( \frac{\partial f}{\partial x} \hat{i} + \frac{\partial f}{\partial y} \hat{j}\right) \cdot \left(\frac{dx}{dt}\hat{i} + \frac{dy}{dt \hat{j}}\right) &= 0\\
   \nabla f(x,y) \cdot \vec{r'}(t) &=0
\end{align*}

\item When is the dot product zero? The vectors are \textbf{orthogonal}. Since $r'(t)$ gives me the slope of the tangent vector to the level curve, the gradient must give us the \textbf{normal} to any point on a level curve. (Also the steepest climb which can be seen with contour maps).

\item If $\nabla f(x,y)$ is the normal to the level curve $f(x,y) = c$, then $\nabla f(x,y,z)$ is the normal to the level surface $f(x,y,z)=c$.

\item So in order to find the normal to a surface, we can find the normal to a level curve just one dimension higher. For example if we have the 2-D curve $x^2+y^2 = 16$ and we want the normal to this curve we need to take the gradient of a surface, i.e. $f(x,y) = x^2-y^2$.



\end{enumerate}




\section{Examples}

\begin{problem}
    Find the tangent plane and the normal line to surface $xyz = -4$ at the point $P(2,-1,2)$
\end{problem}

\newpage

\begin{problem}
    Find the tangent plane and the normal line to the surface $xz^2+yx^2+y^2-2x+3y+6 = 0$ at the point $P(-2,1,3)$.
\end{problem}

\vspace{6 cm}


\begin{problem}
    Find the tangent plane and the normal line to the surface $z = \ln{(x/y)}$ at the point $P(e,1,1)$.
\end{problem}

\vspace{6 cm}


\begin{problem}
    Find the tangent plane and the normal line to the surface $z = \tan^{-1}(y/x)$ at the point $P(1,1,\pi/4))$.
\end{problem}



\newpage

\section{Scratch Work}

\end{document} 